\documentclass[10pt,letterpaper]{article}

%----------------------------------------------------------------------------------------
%	PACKAGES AND THEMES
%----------------------------------------------------------------------------------------
\usepackage[left=0.5in,top=0.5in,right=0.5in,bottom=0.5in]{geometry}
\usepackage{enumitem}
\usepackage{hyperref}
\usepackage{titlesec}
\usepackage{xcolor}
\usepackage{fontawesome5}

% Define links style
\hypersetup{
    colorlinks=true,
    linkcolor=blue,
    filecolor=magenta,
    urlcolor=teal,
}

% Custom section formatting
\titleformat{\section}{
  \vspace{-5pt}\scshape\raggedright\large\bfseries
}{}{0em}{}[\color{black}\titlerule \vspace{-2pt}]

%----------------------------------------------------------------------------------------
%	CUSTOM COMMANDS
%----------------------------------------------------------------------------------------
\newcommand{\resumeItem}[1]{
  \item\small{
    {#1 \vspace{-2pt}}
  }
}

\newcommand{\resumeSubheading}[4]{
  \vspace{-1pt}\item
    \begin{tabular*}{0.99\textwidth}[t]{l@{\extracolsep{\fill}}r}
      \textbf{#1} & \textbf{#2} \\
      \textit{\small#3} & \textit{\small #4} \\
    \end{tabular*}\vspace{-5pt}
}

\newcommand{\resumeSubItem}[2]{\resumeItem{#1}{#2}\vspace{-4pt}}

\renewcommand{\labelitemii}{$\circ$}

\newcommand{\resumeSubHeadingListStart}{\begin{itemize}[leftmargin=*, label={}]}
\newcommand{\resumeSubHeadingListEnd}{\end{itemize}}
\newcommand{\resumeItemListStart}{\begin{itemize}}
\newcommand{\resumeItemListEnd}{\end{itemize}\vspace{-5pt}}

%----------------------------------------------------------------------------------------
%	BEGIN DOCUMENT
%----------------------------------------------------------------------------------------
\begin{document}

%----------HEADING-----------------
\begin{center}
    {\Huge \scshape Sambit Mishra} \\ \vspace{2pt}
    \small PhD Student, Electrical \& Computer Engineering $|$ University of Southern California \\ \vspace{2pt}
    \href{mailto:sambitmi@usc.edu}{\faEnvelope \ sambitmi@usc.edu} $|$ \href{https://sambitmishra.in}{\faGlobe \ sambitmishra.in} $|$ \href{https://scholar.google.com/citations?user=kyCSMKUAAAAJ}{\faGraduationCap \ Google Scholar} $|$ \href{https://orcid.org/0000-0002-7736-7164}{\faOrcid \ ORCID}
\end{center}

%-----------RESEARCH INTERESTS-----------------
\section{Research Interests}
\small{
Causal inference and probabilistic graphical models, with emphasis on identifiability theory and scalable continuous-optimization methods for causal discovery. Interested in efficient discovery algorithms for identifiable DAGs, variance bounds for non-/partially-identifiable cases, and high-dimensional, mixed-distribution data.
}
\vspace{-8pt}

%-----------EDUCATION-----------------
\section{Education}
  \resumeSubHeadingListStart
    \resumeSubheading
      {University of Southern California (USC)}{Los Angeles, CA, USA}
      {Ph.D. in Electrical and Computer Engineering; Advisor: Prof. Urbashi Mitra}{Aug 2025 -- Present}
      \resumeItem{\textbf{GPA: 4.00/4.00}}
      \resumeItem{\textbf{Key Coursework:} Probability Theory, Linear Algebra, Supervised Machine Learning, Inference \& Estimation Theory.}

    \vspace{10pt}

    \resumeSubheading
      {Indian Institute of Technology (IIT) Bhubaneswar}{Bhubaneswar, Odisha, India}
      {B.Tech. (Hons.) in Electronics and Communication Engineering; Advisor: Dr. Soumya P. Dash}{2021 -- 2025}
      \resumeItem{\textbf{GPA: 9.14/10.00}}
  \resumeSubHeadingListEnd

%-----------RESEARCH EXPERIENCE-----------------
\section{Research Experience}
  \resumeSubHeadingListStart

    \resumeSubheading
  {Graduate Research Assistant}{Los Angeles, CA, USA}
  {Ming Hsieh Department of Electrical and Computer Engineering, USC}{Aug 2025 -- Present}
  \begin{itemize}[leftmargin=*, noitemsep]
    \resumeItem{Developing continuous optimization methods for scalable DAG estimation using differentiable constraints to reduce iterations and enable large-scale batching.}
    \resumeItem{Designed a soft-intervention selection framework to improve causal discovery efficiency and lower experimental costs.}
    \resumeItem{Studying identifiability of causal graphs under finite samples and latent variable settings.}
    \resumeItem{Built GPU-accelerated structure learning pipelines in PyTorch using Adam and L-BFGS for non-convex optimization.}
  \end{itemize}

    \resumeSubheading
      {Undergraduate Researcher}{Bhubaneswar, Odisha, India}
      {School of Electrical and Computer Sciences, IIT Bhubaneswar}{Aug 2024 -- May 2025}
      \begin{itemize}[leftmargin=*, noitemsep]
        \resumeItem{Developed SER-optimized modulation schemes for RIS-assisted non-coherent communication systems.}
        \resumeItem{Conducted error analysis for multi-level ASK modulations, optimizing energy efficiency.}
      \end{itemize}

  \resumeSubHeadingListEnd

%-----------PUBLICATIONS (HIGHLIGHTED)-----------------
\section{Publications}
\small \begin{enumerate}[leftmargin=*, label={[\arabic*]}]
    \item  C. Peng, \textbf{S. Mishra}, U. Mitra, ``Learning to Intervene: Optimized Soft Intervention Selection for Causal Discovery,'' in \textit{Proc. IEEE Int. Conf. on Acoustics, Speech and Signal Processing (ICASSP)}, Barcelona, Spain, 2026, pp. 6196--6200.

    \item \textbf{S. Mishra}, S. P. Dash and G. C. Alexandropoulos,``SER-Optimized Multi-Level ASK Modulations for RIS-Assisted Communications With Energy- and Sign-Based Noncoherent Reception," in \emph{IEEE Transactions on Green Communications and Networking}, vol. 10, pp. 1433-1445, 2026, doi: 10.1109/TGCN.2025.3633182.

    \item \textbf{S. Mishra} and S. P. Dash,``Error Analysis With Optimal Receiver and Multi-Level ASK for RIS-Assisted Noncoherent Wireless System," in \emph{IEEE Wireless Communications Letters}, vol. 15, pp. 300-304, 2026, doi: 10.1109/LWC.2025.3624154.

\end{enumerate}

\section{Awards \& Academic Distinctions}
\begin{itemize}[leftmargin=*, noitemsep]
  \item B.Tech. Project Proficiency Award --- Best B.Tech Thesis, ECE Department, IIT Bhubaneswar (2024--25)
  \item Secured All India Rank 4000 among $\sim$1{,}000{,}000 candidates in JEE Advanced 2021, India
  \item Awarded National Talent Search Scholarship by achieving a rank among the top 2000 students among $\sim$1{,}200{,}000 candidates in National Talent Search Examination (NTSE) 2019, Government of India
\end{itemize}

\section{Technical Skills}
 \begin{itemize}[leftmargin=*, noitemsep]
    \item \textbf{Causal Inference:} Causal Discovery, Bayesian Networks, Structural Equation Models
    \item \textbf{Optimization:} Gradient-based solvers (Adam, RMSProp, L-BFGS), Convex Optimization, Bilevel Optimization.
    \item \textbf{Programming \& Tools:} Python (PyTorch, NumPy, SciPy), C++, CUDA, MATLAB, Git, LaTeX.
 \end{itemize}

%-----------INDUSTRY EXPERIENCE (COMPRESSED)-----------------
\section{Industry Experience}
  \resumeSubHeadingListStart
    \resumeSubheading
      {NVIDIA Corporation}{Bengaluru, Karnataka, India}
      {ASIC Engineering Intern, CHI VIP Team}{May 2024 -- July 2024}
      \vspace{8pt}
      \newline
      Verified CHI protocol compliance for CPU/GPU IPs using SystemVerilog and Perl automation; ensured interoperability between internal and external verification IPs.
  \resumeSubHeadingListEnd

\end{document}
